\section{Application to a formula of Le--Murakami and Furusho type}

In this section, we assume that $R$ is a field and consider 
$R\la X,Y\ra$ as a subalgebra of the ring of non-commutative
formal power series $R\lala X,Y\rara$, where a standard
comultiplication $\Delta$ is defined by setting
$\Delta(a)=1\otimes a+a\otimes 1$ for $a\in\{X,Y\}$.  
An element $J\in R\lala X,Y\rara$ is called group-like if
it has constant term 1 and satisfies $\Delta(J)=J\otimes J$. 
There are many group-like elements: 
for example, the subgroup multiplicatively generated by $\exp(X)$ and $\exp(Y)$
in $R\lala X,Y\rara^\times$ consists of group-like elements and forms a free group of rank 2. 

\begin{Theorem}[Le--Murakami, Furusho type formula]
\label{LeMuFu}
Let $J\in R\lala X,Y\rara$ be a group-like element in the form
$$
J=\sum_{\bk\in\N_0^{(\infty)}} c_\bk w_\bk,
$$
and write $c_X$ for the coefficient $c_{(;1)}$ of $X$ in $J$. 
Then,
$$
c_{(k_1,\dots,k_d;k_\infty)}= \hspace{-0.2cm} \sum_{\substack{s,t\ge 0 \\ s+t=k_\infty}}
(-1)^s
\frac{{(c_X)}^t}{t!}
\hspace{-0.2cm} \sum_{\substack{ s_1,\dots,s_d\ge 0 \\ s=s_1+\dots+s_d}}
\binom{k_1+s_1}{k_1}\cdots \binom{k_d+s_d}{k_d}
c_{(k_1+s_1,\dots,k_d+s_d;0)}.
$$
\end{Theorem}


We first prove an elementary identity that will be used for the proof of the above formula.

\begin{Lemma} \label{lem2}
Let $\bkappa=(k_1,\dots,k_d)\in\N_0^d$ and $\bs=(s_1,\dots,s_d)\in \Z^d$
satisfy $s=s_1+\cdots+s_d\ge 0$ and $k_i+s_i\ge 0$ $(i=1,\dots,d)$.
Then,  we have
$$
\sum_{\btau\in\N_0^d} 
\lp\SS^{(\btau;0)},w_{(\bkappa+\bs;0)}\rp
\cdot \lp\MM^{(\btau;0)}, w_{(\bkappa;s)}\rp
=(-1)^s
\binom{k_1+s_1}{k_1}\cdots \binom{k_d+s_d}{k_d}.
$$
\end{Lemma}

\begin{proof}
We shall compute the left-hand side explicitly
as the sum over 
$\btau\in\N_0^d$ satisfying
$\sum_{i=1}^d t_i=\sum_{i=1}^d (k_i+s_i)$
with
$$
\la\SS^{(\btau;0)},w_{(\bkappa+\bs;0)}\ra=\binom{(\bkappa+\bs;0)}{(\btau;0)}
=\binom{k_1+s_1}{t_1} 
\cdots
\binom{\sum_{i=1}^{d-1}(k_i+s_i)-\sum_{i=1}^{d-2}t_i}{t_{d-1}}
\binom{t_d}{t_d}
$$
by Lemma \ref{dS-mono} and with
\small $$
\la\MM^{(\btau;0)}, w_{(\bkappa;s)}\ra
=\magnus{(\btau;0)}{(\bkappa;s)}
=(-1)^{s+\sum_{i=1}^{d-1}(d-i)(t_i-k_i)}
\binom{t_1}{t_1-k_1}
\cdots
\binom{t_{d-1}}{\sum_{i=1}^{d-1}(t_i-k_i)}
\binom{t_d}{s}
$$ \normalsize
by \eqref{Mag-Mono} and
$s=\sum_{i=1}^{d}(t_i-k_i)$. 
Note that,
since $\binom{(\bkappa+\bs;0)}{(\btau;0)}\magnus{(\btau;0)}{(\bkappa;s)}\ne 0$ 
only when all entries of $\btau=(t_1,\dots,t_d)$ are nonnegative
and $t_1+\cdots+t_d=\sum_{i=1}^d (k_i+s_i)$ (a constant),
the above sum
can be taken over the tuples $(t_1,\dots,t_{d-1})\in\Z^{d-1}$
with entries running as independent integers.
Then, the partial summation involved with the last variable 
$t_{d-1}$ may be factored out as
\small \begin{align*}
&\sum_{t_{d-1}}(-1)^{t_{d-1}}
\binom{\sum_{i=1}^{d-1}(k_i+s_i)-\sum_{i=1}^{d-2}t_i}{t_{d-1}}
\binom{t_{d-1}}{\sum_{i=1}^{d-1}(t_i-k_i)}
\binom{t_d}{s} \\
&=\sum_{t_{d-1}}(-1)^{t_{d-1}}
\binom{\sum_{i=1}^{d-1}(k_i+s_i)-\sum_{i=1}^{d-2}t_i}{ \sum_{i=1}^{d-1}k_i-\sum_{i=1}^{d-2}t_i }
\binom{\sum_{i=1}^{d-1} s_i}{\sum_{i=1}^{d-1}(t_i-k_i)}
\binom{\sum_{i=1}^d(k_i+s_i)-\sum_{i=1}^{d-1}t_i}{\sum_{i=1}^{d} s_i} \\
&=\binom{\sum_{i=1}^{d-1}(k_i+s_i)-\sum_{i=1}^{d-2}t_i}{\sum_{i=1}^{d-1} s_i}
(-1)^{\sum_{i=1}^{d-1}k_i -\sum_{i=1}^{d-2} t_i}
\binom{k_d+s_d}{s_d},
\end{align*} \normalsize
where \cite[(5.21)]{GKP} is applied for the first equality and 
\cite[(5.24)]{GKP} for the second.
After factoring out the constant $\binom{k_d+s_d}{s_d}$ and repeating
the similar process with the other variables $t_{d-2},\dots, t_1$ consecutively,
we eventually obtain the asserted formula.
Below in Note \ref{Quest4.3}, we also provide an alternative proof of the lemma 
free from intricate use of 
\cite[(5.21),(5.24)]{GKP}.
\end{proof}

\begin{proof}[Proof of Theorem \ref{LeMuFu}]
We argue in the beautiful framework exploited in Reutenauer's book \cite[1.5]{Reu}
using the complete tensor product
$$
\mathscr{A}=R\lala X,Y\rara \bar\otimes R\lala X,Y\rara
$$
equipped with a product induced from 
the shuffle product (resp. the concatenation product)
on the left (resp. right) of $\bar\otimes$. 
Recall that the ring of
$R$-linear endomorphisms $\mathrm{End}_RR\lala X,Y\rara$
can be embedded into $\mathscr{A}$ by 
$f\mapsto \sum_{w\in W} w\otimes f(w)$, and that the product of 
$\mathscr{A}$ restricts to the convolution product of $\mathrm{End}_RR\lala X,Y\rara$
defined by $f\ast g:=\mathrm{conc}\circ(f\otimes g)\circ \Delta$
(where ``$\mathrm{conc}$'' means concatenation of left and right sides of $\otimes$).
Note that, for $f\in\mathrm{End}_R R\lala X,Y\rara$ and  $J\in R\lala X,Y\rara$,
we have $f(J)=\sum_{w\in W} \lp w,J\rp f(w)$.

Since, by Corollary \ref{coro}, every word $w$ can be written as 
$\sum_{\bt\in \N_0^{(\infty)}}
\lp \SS^{(\bt)},w \rp \MM^{(\bt)}$, 
the element of $\mathscr{A}$ corresponding to 
the identity $\mathrm{id}\in \mathrm{End}_RR\lala X,Y\rara$ is:
\begin{align*}
\sum_{w\in W} w\otimes w &=
\sum_w w\otimes \sum_\bt \lp \SS^{(\bt)},w \rp\MM^{(\bt)} 
=
\sum_\bt (\sum_w  \lp\SS^{(\bt)},w\rp w) \otimes \MM^{(\bt)} \\
&= 
\sum_\bt \SS^{(\bt)} \otimes \MM^{(\bt)} \\
&=
\left(\sum_{d=0}^\infty 
\sum_{\btau\in\N_0^d} \SS^{(\btau;0)}\otimes \MM^{(\btau;0)}\right)
\cdot
\left(\sum_{t=0}^\infty X^t\otimes X^t\right),
\end{align*}
where we have used
$\SS^{(\bt)}=\SS^{(\btau;t)}=\SS^{(\btau;0)}\sha X^t$ and 
$\MM^{(\bt)}=\MM^{(\btau;t)}=\MM^{(\btau;0)}\cdot X^t$.
Observing that both factors of the last expression above correspond to specific 
$R$-linear endomorphisms,
we can apply $\mathrm{id}$ to $J$ as the convolution product of them and find from
$\Delta(J)=J\otimes J$ that
\begin{equation} \label{J-expansion}
J=\mathrm{id}(J)=
\left(\sum_{d=0}^\infty 
\sum_{\btau\in\N_0^d} \lp \SS^{(\btau;0)},J \rp \MM^{(\btau;0)}\right)
\left(\sum_{t=0}^\infty \frac{(c_X)^t}{t !} X^t\right).
\end{equation}
Note here that the pairing of $J$ with 
$X^t=X^{\sha\, t}/t!$ is equal to $(c_X)^t/t !$,
as easily seen from the fact that the specialization 
$J(X,0)\in R\lala X\rara$ at $Y=0$ is a group like
element $\exp(c_X\cdot X)$.
To complete the proof of Theorem \ref{LeMuFu},
given a fixed $\bk=(\bkappa;k_\infty)=(k_1,\dots,k_d; k_\infty)\in\N_0^{(\infty)}$
and $0\le s\le k_\infty$, 
we compute
the coefficient of~$w_{(\kappa;s)}=X^{k_1}Y\cdots X^{k_d}YX^s$ in the expansion 
of the first factor of the above right-hand side as follows:
\begin{align*}
&\sum_{d=0}^\infty 
\sum_{\btau\in\N_0^d} 
\lp \SS^{(\btau;0)},J \rp 
\lp \MM^{(\btau;0)},w_{(\bkappa;s)} \rp 
=\left\lp\sum_{d=0}^\infty 
\sum_{\btau\in\N_0^d} \lp\SS^{(\btau;0)},J \rp \MM^{(\btau;0)},
w_{(\bkappa;s)} \right\rp \\
=&\left\lp\sum_{d=0}^\infty 
\sum_{\btau\in\N_0^d} 
\biggl\lp\SS^{(\btau;0)},\sum_{\bu\in \N_0^{(\infty)}}(J,w_\bu)w_\bu\biggr\rp
 \MM^{(\btau;0)},
w_{(\bkappa;s)}\right\rp \\
=&\sum_{\bu} \lp J,w_\bu \rp \sum_{d=0}^\infty 
\sum_{\btau\in\N_0^d} \lp \SS^{(\btau;0)},w_\bu \rp
\lp \MM^{(\btau;0)}, w_{(\bkappa;s)} \rp .
\end{align*}
But since $\lp\SS^{(\btau;0)},w_\bu\rp \lp \MM^{(\btau;0)}, w_{(\bkappa;s)} \rp$ 
survives only when 
$\dep(\btau;0)=\dep(\bkappa;s)=\dep(\bu)$ and $|(\btau;0)|=|(\bkappa;s)|=|\bu|$,
the summation $\sum_\bu$ in the last expression above occurs only for 
those $\bu$ of the form 
$(\bkappa+\bs;0)\in \N_0^{(\infty)}$ 
with $\bs=(s_1,\dots,s_d)\in\Z^d$, 
$s=s_1+\cdots+s_d\ge 0$
(cf. also Remark \ref{rem2.6}).  
Then, it follows from Lemma \ref{lem2} that the last expression above is equal to 
$$
\sum_{d=0}^\infty \sum_{\substack{\bs\in\N_0^d \\ |(\bs;0)|=s}} 
\lp J, w_{(\bkappa+\bs;0)} \rp
(-1)^s
\binom{k_1+s_1}{k_1}\cdots \binom{k_d+s_d}{k_d}.
$$
(Note: the prescribed condition $\bs\in\Z^d$ has been replaced 
with $\bs\in\N_0^d$ because of the a posteriori survivals of binomial factors.)
From this and \eqref{J-expansion} together with
$\lp J, w_{(\bkappa+\bs;0)} \rp=c_{(k_1+s_1,\dots,k_d+s_d;0)}$, 
we conclude the assertion.
\end{proof}

\begin{Note}[{\it Alternative proof of Lemma} \ref{lem2}] \label{Quest4.3}
In the right-hand side of Lemma \ref{lem2}, the quantity
$\binom{k_1+s_1}{k_1}\cdots \binom{k_d+s_d}{k_d}$ appearing there can also be interpreted as
the pairing
$\lp w_{(k_1,\dots,k_d;0)}\sha X^s, w_{(k_1+s_1,\dots,k_d+s_d;0)} \rp$.
Therefore, the assertion of this lemma is equivalent to the identity
\begin{equation}
\sum_{\btau\in\N_0^d} 
\lp\SS^{(\btau;0)},w_{(\bkappa+\bs;0)} \rp \cdot \lp\MM^{(\btau;0)}, w_{(\bkappa;0)}\!\cdot\!\! X^s\rp
=(-1)^s 
\lp w_{(\bkappa;0)}\sha X^s, w_{(\bkappa+\bs;0)} \rp
\label{alterLemma}
\end{equation}
for $\bkappa=(k_1,\dots,k_d)\in\N_0^d$, $\bs=(s_1,\dots,s_d)\in\Z^d$ 
satisfying $s=s_1+\cdots+s_d\ge 0$ and $\bkappa+\bs\in\N_0^d$.
We now give an alternative proof for it using the Magnus/demi-shuffle duality.
First, by Corollary \ref{coro}, we have 
$w_{(\bkappa;0)}=\sum_\bbr\la \MM^{(\bbr)},w_{(\bkappa;0)}\ra \SS^{(\bbr)}$ 
and $w_{(\bkappa+\bs;0)}=\sum_\bt\la \SS^{(\bt)},w_{(\bkappa+\bs;0)}\ra \MM^{(\bt)}$
so that the right-hand side of \eqref{alterLemma} can be written as
\begin{align}
\label{RHS4.2}
&(-1)^s 
\lp w_{(\bkappa;0)}\sha X^s, w_{(\bkappa+\bs;0)} \rp \\
&=(-1)^s \sum_{\bbr,\bt\in\N_0^{(\infty)}}
\lp \SS^{(\bbr)}\sha X^s, \MM^{(\bt)} \rp
\lp \MM^{(\bbr)} ,w_{(\bkappa;0)} \rp
\lp \SS^{(\bt)} ,w_{(\bkappa+\bs;0)}\rp
\notag \\
&=(-1)^s \sum_{\brho\in\N_0^d} 
\lp \MM^{(\brho;0)} ,w_{(\bkappa;0)}\rp
\lp  \SS^{(\brho;s)}, w_{(\bkappa+\bs;0)} \rp.
\notag
\end{align}
Here in the second equality, we use the fact that
$\lp  \MM^{(\bbr)}, w_{(\bkappa;0)} \rp$ 
survives only if $\bbr=(\brho;0)\in\N_0^{(\infty)}$ for some $\brho\in\N_0^d$
and then apply the duality (Theorem \ref{orthogonality})
to $\lp \SS^{(\bbr)}\sha X^s, \MM^{(\bt)} \rp$
with $\SS^{(\brho;0)}\sha X^s=\SS^{(\brho;s)}$ 
(cf. Definitions \ref{defMagnus} and \ref{demi-shuffle}).

On the other hand, in the left-hand side of  \eqref{alterLemma},
one observes that nontrivial terms of the summation arise
only from those $\btau=(\tau_1,\dots,\tau_d)\in\N_0^d$ with
$\sum_{i=1}^d \tau_i=s+\sum_{i=1}^d \kappa_i$ (a constant). 
But the last binomial factor in \eqref{MagnusBino} for
$ \lp\MM^{(\btau;0)}, w_{(\bkappa;0)}\!\cdot\!\! X^s\rp
=\magnus{(\tau_1,\dots,\tau_d;0)}{(\kappa_1,\dots,\kappa_d;s)}
$
equals $\binom{\tau_d}{s}$, which is non-zero only if
$\tau_d\ge s$.
Therefore, 
the summation~$\sum_\btau$ may be replaced by 
$\sum_\brho$ with $\brho=\btau-(\mathbf{0},s)$
in $\N_0^d$ (where $\mathbf{0}\in\N_0^{d-1}$ is the zero vector).
Thus,
the left-hand side of \eqref{alterLemma} can be written as
\begin{align}
\label{LHS4.2}
&\sum_{\btau\in\N_0^d} 
\lp\SS^{(\btau;0)},w_{(\bkappa+\bs;0)} \rp \cdot 
\lp\MM^{(\btau;0)}, w_{(\bkappa;0)}\!\cdot\!\! X^s\rp \\
&=
\sum_{\brho\in\N_0^d} 
\lp\SS^{(\brho+(\mathbf{0},s);0)},w_{(\bkappa+\bs;0)} \rp \cdot 
\lp\MM^{(\brho+(\mathbf{0},s);0)}, w_{(\bkappa;s)} \rp.
\notag
\end{align}
Comparing summands of the above \eqref{RHS4.2} and \eqref{LHS4.2} for 
individual $\brho\in\N_0^d$ in view of coefficients of 
monomial expansions of demi-shuffle/Magnus polynomials 
(Lemma~\ref{dS-mono} and~\eqref{Mag-Mono}), we reduce 
the formula \eqref{alterLemma} to the following elementary identity
for $\bkappa=(k_i)$, $\brho=(r_i)\in\N_0^d$ and $\bs=(s_i)\in\Z^d$
satisfying
$\sum_{i=1}^d k_i=\sum_{i=1}^d r_i$, $\bs+\bkappa\in \N_0^d$
and $s:=\sum_{i=1}^d s_i\ge 0$:
\begin{equation}
\binom{(\bkappa+\bs;0)}{(\brho;s)}
\magnus{(\brho;0)}{(\bkappa;0)}
=
(-1)^s
\binom{(\bkappa+\bs;0)}{(\brho+(\mathbf{0},s);0)}
\magnus{(\brho+(\mathbf{0},s);0)}{(\bkappa,s)}.
\end{equation} 
This is an immediate consequence of definitions
of these symbols $\{^*_* \}$, $(^*_*)$.
(Observe that only difference between the corresponding symbols
occurs from the last binomial coefficient in  \eqref{ArrayBino}
and \eqref{MagnusBino}.)
\qed
\end{Note}
